\section{Conclusions}
We built a packet-to-flow feature extraction pipeline with a sequential
baseline, two OpenMP strategies, a CUDA pipeline with three reduction kernels
and an MPI4Py driver, and verified that all of them produce identical flow
records. On synthetic traffic the GPU version is up to \sizeChunkSp\x{} faster
than the baseline and 8 threads up to \skewPartSpZero\x; on a real backbone
trace the gains are 4.3\x{} and \realPartSp\x.

The main lesson is that grouping performance depends on the shape of the
traffic more than on the programming model. Many flows make the problem
memory-bound; partitioning by flow hash then helps even with one thread. A
dominant flow breaks every scheme that gives a flow to a single worker
(a partition, a GPU thread) and makes a shared table slower than one thread;
schemes that split the packets evenly and combine partial results (the merge
on the processor, the chunked kernel on the GPU) are robust to it. Real traffic
shows both problems at once, and its many single-packet flows make the output
transfer the largest GPU cost. Parallel resources beyond the 4 physical cores
of the test machine gave no benefit, for threads and for processes alike.

Future work: choose the strategy at run time from a sample of the batch; run
the MPI driver on several nodes, where memory is not shared; process
captures larger than GPU memory in pieces using the merge operation; and add
bidirectional flows and flow timeouts as flow exporters use.
